%% kapa: exam=1 round=22 year=2011 period=1 form=A subject=economics subject_ko=경제원론 seq=17 number=57
%% source: Q-Net 감정평가사 기출문제 (cst00309) artlSeq=5036790 file=extracted/1008/5036790/2011년 제22회 감정평가사 1차시험 기출문제/1교시.hwp
%% conversion: kordoc 4.24.51 -> tools/gen1.py (auto); review=proofread-sonnet; flags=none
\begin{question}{57}
甲과 乙이 총 금액 10만원을 나누어 갖는 2인 비협조게임에서 규칙은 보기와 같다. 다음 전략 중 내쉬균형에 해당하는 것은?
\begin{bogi}
- 甲과 乙이 각각 10만원 미만에서 만원단위로 자기가 원하는 금액을 동시에 제시한다. \par
- 甲과 乙이 제시한 금액의 합이 10만원을 초과하고 제시한 금액이 동일할 경우 각각 5만원씩을 받으며, 제시한 금액이 서로 다를 경우 적은 금액을 제시한 사람은 자신이 제시하는 금액을 받고, 더 많은 금액을 제시한 사람은 나머지 금액만을 받는다. \par
- 甲과 乙이 제시한 금액의 합이 10만원 이하일 경우 각자 제시한 금액을 받고 10만원에서 남은 금액이 있으면 폐기된다.
\end{bogi}
\begin{choices}
  \item 甲 2만원, 乙 9만원
  \item 甲 4만원, 乙 6만원
  \item 甲 5만원, 乙 6만원
  \item 甲 8만원, 乙 2만원
  \item 甲 9만원, 乙 1만원
\end{choices}
\end{question}
\answer{3}
